\documentclass[../../main.tex]{subfiles}
\begin{document}

{\bf [解]}
$A=\{20,21,22,23,24\}$，$B=\{25,26,27,28,29\}$，$C=\{80,81,82,83,84\}$，
$D=\{85,86,87,88,89\}$ とする．題意の試行では，$A,B$ のうちから1つ，$C,D$ のうちから1つ数をとり出す．その組は以下の4通りに分けて考えられる．

\begin{center}
\begin{tabular}{c|c|c}
 & $A$ & $B$ \\
\hline
$C$ & ア & イ \\
\hline
$D$ & ウ & エ
\end{tabular}
\end{center}

このうち，アでは必ず $S\ge S'$，エでは必ず $S<S'$ となる．エは$25$通りありあるから，
\begin{align} 
 \begin{cases}
 P_{\text{ア}} = 0 \\
 P_{\text{エ}} = 25
 \end{cases} \label{eq:1}
\end{align}
通りが題意の条件を満たす．

そこで，以下イ，ウについて考える．

\subparagraph{$1.$ イの時}

$B,C$ の数は四捨五入すると各々30,80になるから，$S'=2400$ である．
$S$ を具体的に計算して$S'$との大小を計算すると以下のようになる．ただし$\bigcirc$ は $S<S'$ を，$\times$ は $S\ge S'$ を表す．
すなわち全25通りのうち
\begin{align}
 P_{\text{イ}} = 23 \label{eq:2}
\end{align}
通りが題意の条件を満たす．

\begin{table}[htb]
 \begin{center}
\begin{tabular}{c|c|c|c|c|c}
$C\backslash B$ & 25 & 26 & 27 & 28 & 29 \\
\hline
80 & $\bigcirc$ & $\bigcirc$ & $\bigcirc$ & $\bigcirc$ & $\bigcirc$ \\
\hline
81 & $\bigcirc$ & $\bigcirc$ & $\bigcirc$ & $\bigcirc$ & $\bigcirc$ \\
\hline
82 & $\bigcirc$ & $\bigcirc$ & $\bigcirc$ & $\bigcirc$ & $\bigcirc$ \\
\hline
83 & $\bigcirc$ & $\bigcirc$ & $\bigcirc$ & $\bigcirc$ & $\times$ \\
\hline
84 & $\bigcirc$ & $\bigcirc$ & $\bigcirc$ & $\bigcirc$ & $\times$
\end{tabular}
\end{center}
\end{table}


\subparagraph{$2^\circ$ ウの時}

$A,D$ の数は，四捨五入すると各々20,90になるから，$S'=1800$ である．$S$ について同様に計算すると以下の表を得る，

\begin{table}[htb]
 \begin{center}
\begin{tabular}{c|c|c|c|c|c}
$D\backslash A$ & 20 & 21 & 22 & 23 & 24 \\
\hline
85 & $\bigcirc$ & $\bigcirc$ & $\times$ & $\times$ & $\times$ \\
\hline
86 & $\bigcirc$ & $\times$ & $\times$ & $\times$ & $\times$ \\
\hline
87 & $\bigcirc$ & $\times$ & $\times$ & $\times$ & $\times$ \\
\hline
88 & $\bigcirc$ & $\times$ & $\times$ & $\times$ & $\times$ \\
\hline
89 & $\bigcirc$ & $\times$ & $\times$ & $\times$ & $\times$
\end{tabular}
\end{center}
\end{table}
従って題意の条件を満たすのは
\begin{align}
 P_{\text{ウ}} = 6 \label{eq:3}
\end{align}
通りである．

以上で全ての場合分けは尽くされた．
\cref{eq:1,eq:2,eq:3}より，全てのえらび方 $10^2=100$ 通りのうち $S<S'$ となるのは，
\begin{align}
P_{\text{ア}}+P_{\text{エ}}+P_{\text{イ}}+P_{\text{ウ}} = 25+23+6=54 \text{（通り）}
\end{align}
である．よって求める確率は
\begin{align}
\dfrac{54}{100}=\dfrac{27}{50} 
\end{align}
である．

\newpage

\end{document}
